% v24. The four-ray isolated-vertex lemma shortens the v18 proof.
% The proof is unbounded; all integer-gap arguments are explicit.
\subsection{三高分母の正符号グラフは完全マッチングである}
\label{mh18:section}

命題~\ref{sm14:four-sign-types} の四つの候補から，整数性によって二つの道型を除き，
正の weights によって三角形型を除く．これにより，以前は仮定していた完全マッチング性が
全ての正の三高分母模型について成り立つ．

\begin{theorem}[三高分母の完全マッチング定理]
\label{mh18:matching}
$q_0=1$，$q_1,q_2,q_3>1$ とし，整数 $p_{ij}$ が行和と pair products
\[
 \sum_{j\ne i}p_{ij}=-q_i,\qquad p_{ij}p_{ji}=q_iq_j-1
\]
を満たすとする．$p_{ij}>0$ のとき無向辺 $ij$ を置くグラフを $G^+$ とする．
このとき次が成り立つ．
\begin{enumerate}
\item $G^+$ は完全マッチング，または低 ray $0$ を含む三角形と孤立頂点の和である．
\item 正の basin weights が存在することと，$G^+$ が完全マッチングであることは同値である．
\end{enumerate}
特に正の模型の低 row は，ちょうど二つの負の entry と一つの正の entry を持つ．
二つの負の low-row entries の相手同士が，もう一つの正の辺を作る．
分母の上限，整数 clutching，対向 $D$ の相殺，node 条件は仮定しない．
\end{theorem}

\begin{proof}
命題~\ref{sm14:four-sign-types} により，除くべき型は，$0$ を端点に持つ $P_3$ と孤立頂点，
$0$ を内部頂点に持つ $P_4$，および $0$ を含む三角形である．

\paragraph{$0$ が端点である $P_3$ の排除．}
正の辺を $03,13$ とし，正整数 $s,t,S,T,U$ により
\[
 (p_{01},p_{02},p_{03})=(-s,-t,s+t-1),\quad
 (p_{10},p_{20},p_{30})=(-S,-T,U)
\]
と書く．$u=s+t-1$ とすると，$q_1=sS+1,q_2=tT+1,q_3=uU+1$ である．
さらに
\[
 p_{13}=L>0,\quad p_{31}=M>0,\quad
 (p_{12},p_{21},p_{23},p_{32})=(-H,-K,-N,-O)
\]
と置き，
\[
 a=q_1q_2-1,\quad b=q_2q_3-1,\quad d=q_1q_3-1,
\]
\[
 B=q_1-S>0,\quad C=q_3+U>0,\quad k=q_2-T>0
\]
と略記する．行和と pair equations は
\[
 H=L+B,\quad O=M+C,\quad K+N=k,\quad
 HK=a,\ NO=b,\ LM=d
\]
を与える．従って
\begin{equation}\label{mh18:path-polynomial}
 (kC-b)L^2+\{kd-aC+B(kC-b)\}L+d(kB-a)=0.
\end{equation}
一方，直接計算により
\[
 kC-b=-U(sT-1)-T+1\le0,
\]
\[
 kB-a=-uST-S-T+1<0,\qquad
 kd-aC=-dT-aU+q_3-q_2<0.
\]
最後の不等式は $d=q_1q_3-1>q_3$ と $T\ge1$ による．
式~\eqref{mh18:path-polynomial} は $L>0$ で狭義負となるため矛盾する．

\paragraph{$0$ が内部頂点である $P_4$ の排除．}
正の辺を $01,02,13$ とし，今度は
\[
 (p_{01},p_{02},p_{03})=(s,t,-s-t-1),\quad
 (p_{10},p_{20},p_{30})=(S,T,-U)
\]
と書く．$u=s+t+1$，$q_1=sS+1,q_2=tT+1,q_3=uU+1$ とし，
$L,M,H,K,N,O,a,b,d$ は前段落と同じ符号で定める．
\[
 B=q_1+S,\quad C=q_3-U,\quad k=q_2+T
\]
と置けば，全て正であり，同じ式~\eqref{mh18:path-polynomial} が必要となる．
この場合は
\[
 kC-b=U(sT-1)+T+1>0,
\]
\[
 kB-a=uST+S+T+1>0,\qquad
 kd-aC=dT+aU+q_3-q_2>0.
\]
最後は $a=q_1q_2-1>q_2$，$U\ge1$ による．
式の左辺は $L>0$ で狭義正となり，再び矛盾する．
ここまでには正の weights を用いていない．

\paragraph{三角形に対する正値性の障害．}
三角形と孤立頂点の和は，補題~\ref{fg24:no-isolated} により正の weights を持たない．
この補題は，孤立頂点の負の row から作る二つの整数 gap の和が一になることを用いる．
各 row で非正の gap を選ぶと，gain が一以上の有向閉路が生じるためである．
従って低 ray を含む三角形型も排除される．

\paragraph{完全マッチングの十分性．}
正の二辺は互いに端点を共有しない．各正の辺 $ij$ について
\[
 \frac{p_{ij}}{q_i}<\frac{w_j}{w_i}<\frac{q_j}{p_{ji}}
\]
という開区間は，$p_{ij}p_{ji}=q_iq_j-1<q_iq_j$ により空でない．
二辺について独立に有理数の比を選べば，正の有理 weights が得られる．
負の entries の basin inequalities は自動的に成立する．これで同値性も従う．
\end{proof}

\begin{remark}[正値性の仮定と適用範囲]
\label{mh18:positivity-essential}
正値性を落とすと三角形型は実際に存在する．例えば
\[
 q=(1,5,19,31),\qquad
 P=\begin{pmatrix}
 0&1&1&-3\\
 4&0&2&-11\\
 18&47&0&-84\\
 -10&-14&-7&0
 \end{pmatrix}
\]
は全ての行和・pair products を満たす．しかし正の閉路 $0\to2\to1\to0$ の gain は
$188/95>1$ であり，正の weights は存在しない．
完全マッチング定理は，三つの高 $D$ が同時に非負でないことをまだ主張しない．
一方，\ref{sh14:matching-setup} の Schur 表示と
定理~\ref{sm14:beta-negative} の $\beta\le-1$ は，今後，全ての正の三高分母模型に適用できる．
\end{remark}
